\subsection{Особенности адресации (Буфер драйвера)}

Важно понимать, как драйвер светодиодов обрабатывает команды. Модуль \hyperlink{term:ledbar}{Ledbar} работает с внутренним \hyperlink{term:buffer}{буфером памяти}.

Когда вы создаёте объект \texttt{Ledbar.new(N)}, система выделяет память под \texttt{N} светодиодов.

\begin{warnbox}[Критическая ошибка]
Стандартная библиотека требует инициализации буфера размером \textbf{25 элементов}, даже если физически на плате распаяно всего 4 светодиода.
\begin{itemize}
    \item \textbf{Правильно:} \texttt{Ledbar.new(25)}
    \item \textbf{Неправильно:} \texttt{Ledbar.new(4)} — в этом случае светодиоды могут работать некорректно или не включиться вовсе из-за особенностей прошивки драйвера.
\end{itemize}
\end{warnbox}

Типовой шаблон инициализации выглядит так:

\begin{minted}{lua}
local unpack = table.unpack
local ledNumber = 25                  -- размерность буфера (всегда 25!)
local leds = Ledbar.new(ledNumber)    -- объект управления лентой

local function changeColor(col)
  for i=0, ledNumber-1 do             -- проходим по всем индексам 0..24
    leds:set(i, unpack(col))          -- записываем RGB в буфер
  end
end

changeColor({0,1,0})                  -- включаем зелёный
\end{minted}
